\documentclass[11pt]{article}

\usepackage[a4paper,margin=1in]{geometry}
\usepackage[T1]{fontenc}
\usepackage[utf8]{inputenc}
\usepackage{lmodern}
\usepackage{microtype}
\usepackage{amsmath,amssymb,amsthm,mathtools}
\usepackage{array}
\usepackage{enumitem}
\usepackage{xurl}
\usepackage[colorlinks=true,linkcolor=blue,citecolor=blue,urlcolor=blue]{hyperref}

\theoremstyle{plain}
\newtheorem{theorem}{Theorem}[section]
\newtheorem{proposition}[theorem]{Proposition}
\newtheorem{lemma}[theorem]{Lemma}
\newtheorem{corollary}[theorem]{Corollary}
\theoremstyle{definition}
\newtheorem{definition}[theorem]{Definition}
\newtheorem{hypothesis}[theorem]{Hypothesis}
\theoremstyle{remark}
\newtheorem{remark}[theorem]{Remark}

\DeclareMathOperator{\cyc}{cyc}
\DeclareMathOperator{\rank}{rank}
\DeclareMathOperator{\Card}{Card}
\newcommand{\N}{\mathbb N}
\newcommand{\R}{\mathbb R}
\newcommand{\Sym}{\mathfrak S}
\newcommand{\dd}{\mathrm d}
\newcommand{\one}{\mathbf 1}
\newcommand{\g}{\gamma}
\newcommand{\p}{\pi}
\newcommand{\defp}{\delta_{+}}
\newcommand{\defm}{\delta_{-}}
\newcommand{\Def}{\Delta}
\newcommand{\Env}{\mathcal E}
\newcommand{\LT}{\mathcal L}
\newcommand{\NC}{\mathrm{NC}}
\newcolumntype{L}[1]{>{\raggedright\arraybackslash}p{#1}}

\title{A Bidefect Counting Reduction for Aubrun's Partial-Transpose Moment Bound}
\author{Kieran McShane}
\date{June 7, 2026}

\begin{document}
\maketitle

\begin{abstract}
We isolate a pure combinatorial reduction underlying Aubrun's moment method
for partial transposes of random induced states.  The reduction replaces a
single total-defect count by a two-parameter count indexed by the two oriented
Cayley defects of a Wick permutation.  This bidefect form is the one compatible
with the shifted scalar bases appearing in the analytic estimate.

The main result of this draft is conditional and deliberately explicit: the
bidefect envelope follows from two large-gap Biane adjacent-link inputs, one
nested predecessor-skip collision input, and one Chapuy-type genus recurrence.
The finite gap-three Biane branches and the direct predecessor-skip branch are
separated from these four inputs and accounted for by elementary
incidence-tree reductions.  A new branch-normalized lemma records the
canonical first/last return interval in every genuinely nested predecessor
skip, and the current reduction exposes the nested collision input through
this return interval.  The remaining four inputs are stated in mathematical
form rather than hidden in notation.
\end{abstract}

\tableofcontents

\section{Purpose and Scope}

Aubrun's proof of the partial-transpose moment estimate is a moment-counting
argument over Wick permutations.  The analytic side assigns to each permutation
\(\p\) a power of the matrix dimension and the environment dimension; the
combinatorial side must show that the number of permutations with a given loss
of rank is not too large.  A tempting formulation counts only the total loss.
That formulation is not stable enough for the final scalar estimate, because
the two losses are paid by two different shifted bases.

This article develops the corrected two-parameter formulation.  It is written
as an ordinary mathematical draft.  The formalization status is recorded only
in Appendix~\ref{sec:sync}; the main body is meant to be read without Lean.

\begin{remark}[Honesty of the present version]
The unconditional U-AUB-02 theorem is not closed in the formal development at
the time of this draft.  The theorem below is therefore a reduction theorem:
it says exactly which four combinatorial inputs remain.  It is nevertheless a
usable article skeleton, because those inputs are now localized and no longer
confused with rank-budget bookkeeping or finite checked branches.
\end{remark}

\section{Permutations, Defects, and the Bidefect Envelope}

Fix \(m\geq 0\), put \(n=m+1\), and let
\[
  \g=(0\,1\,\cdots\,m)\in\Sym_n
\]
be the long cycle.  For a permutation \(\sigma\), write \(\cyc(\sigma)\) for
the number of cycles, including fixed points.  The Cayley length is
\[
  |\sigma|=n-\cyc(\sigma).
\]

In the partial-transpose Wick expansion, the surviving index identifications
can be encoded by a permutation \(\p\in\Sym_n\).  This is the only meaning of
``Wick permutation'' in this article.  The sample index contributes
\(2\cyc(\p)\) in the balanced \(s\sim \lambda d^2\) regime, while the two
matrix-index directions contribute \(\cyc(\g\p)\) and \(\cyc(\p\g^{-1})\).
This is the source of the weighted rank in Section~3.

\begin{definition}[Oriented defects]
For \(\p\in\Sym_n\), define the two oriented defects
\[
  \defp(\p) = n+1-\cyc(\p)-\cyc(\g\p),
  \qquad
  \defm(\p) = n+1-\cyc(\p)-\cyc(\p\g^{-1}).
\]
The bidefect class and its cardinality are
\[
  \mathcal B_m(a,b)
  =
  \{\p\in\Sym_n:\defp(\p)=a,\ \defm(\p)=b\},
  \qquad
  N_m(a,b)=|\mathcal B_m(a,b)|.
\]
\end{definition}

The inequalities
\[
  \cyc(\p)+\cyc(\g\p)\leq n+1,
  \qquad
  \cyc(\p)+\cyc(\p\g^{-1})\leq n+1
\]
are just the Cayley triangle inequality written in cycle-count form.  Hence
\(\defp\) and \(\defm\) are non-negative integers.

\begin{definition}[Aubrun bidefect envelope]
Let \(P(Q)\) be the polynomial scalar envelope used in the Wick estimate.  In
the formal development it is denoted \(P(Q)\) with
\[
  Q=2(m+1).
\]
Define
\[
  \Env_m(a,b)
  =
  2^{2(m+1)}
  \binom{m+3}{a}
  \binom{m+1}{b}
  P(2(m+1))^{a+b}.
\]
The bidefect count target is
\[
  N_m(a,b)\leq \Env_m(a,b)
\]
for
\[
  0\leq a\leq m+3,\qquad 0\leq b\leq m+1.
\]
\end{definition}

\begin{remark}
The ranges are part of the theorem.  One can state an all-\((a,b)\) version
only after proving that the out-of-range classes are empty or harmless.  The
current reduction deliberately avoids that extra claim.
\end{remark}

\section{Why the Count Must be Bidefect}

The Wick weight contains three cycle-count contributions: one from the sample
index and two from the left and right matrix indices.  In the balanced regime
the weighted rank has the form
\[
  R(\p)=2\cyc(\p)+\cyc(\g\p)+\cyc(\p\g^{-1}).
\]
The rank budget is
\[
  R(\p)+\defp(\p)+\defm(\p)=2n+2.
\]
Equivalently, if the total defect is
\[
  \Def(\p)=2n+2-R(\p),
\]
then
\[
  \Def(\p)=\defp(\p)+\defm(\p).
\]

The identity explains both the success and the failure of the old route.  A
single-defect count sees the correct rank loss, but it forgets which side of
the scalar estimate pays that loss.  The bidefect envelope keeps the two
losses separate until the final scalar absorption.

\begin{proposition}[Shifted-base absorption, interface form]
Assume the bidefect count target
\[
  N_m(a,b)\leq \Env_m(a,b)
\]
for all \(a,b\) in the stated ranges.  Then the corresponding double sum over
Wick permutations is bounded by the shifted scalar bases used in Aubrun's
moment estimate, rather than by a common maximum base.
\end{proposition}

\begin{proof}
After grouping permutations by \((a,b)\), the count contributes the two
binomial factors
\[
  \binom{m+3}{a},\qquad \binom{m+1}{b}.
\]
The analytic summand has the schematic form
\[
  A^{m+3-a} B^{m+1-b} Q^{a+b},
\]
where \(A\) and \(B\) are the two scalar bases associated with the two matrix
directions and \(Q=P(2(m+1))\) is the polynomial loss budget.  Thus the
\(a\)-sum is bounded by
\[
  \sum_a \binom{m+3}{a} A^{m+3-a}Q^a=(A+Q)^{m+3},
\]
and the \(b\)-sum is bounded by
\[
  \sum_b \binom{m+1}{b} B^{m+1-b}Q^b=(B+Q)^{m+1}.
\]
The important point is that the exponent \(a\) is absorbed on the left side
and \(b\) on the right side; replacing them by \(a+b\) before the scalar step
loses the shifted-base structure needed for the final estimate.
\end{proof}

\section{The One-Sided Genus Route}

The bidefect estimate is supplied from the plus side.  A bidefect class embeds
in its plus-defect class:
\[
  \mathcal B_m(a,b)\subseteq \{\p:\defp(\p)=a\}.
\]
Thus the bidefect problem is reduced to a one-sided count for \(\defp\), with
the missing \(b\)-dependence restored by the binomial scalar envelope.

\begin{definition}[Left-tight permutations]
A permutation \(\p\in\Sym_n\) is left-tight if
\[
  \cyc(\p)+\cyc(\g\p)=n+1.
\]
Equivalently, \(\defp(\p)=0\).
\end{definition}

\begin{proposition}[Genus-zero base]
The zero plus-defect class is the class of left-tight permutations.  Therefore
the genus-zero count is bounded once left-tight permutations inject into
noncrossing partitions:
\[
  \#\{\p:\defp(\p)=0\}
  \leq |\NC(n)|
  \leq 4^n.
\]
\end{proposition}

\begin{proof}
The first statement is the definition of \(\defp\).  The second follows if the
cycle partition of a left-tight permutation is noncrossing and determines the
permutation.  The number of noncrossing partitions on \(n\) points is the
\(n\)-th Catalan number, and the elementary bound \(C_n\leq 4^n\) gives the
displayed estimate.
\end{proof}

For positive plus-defect the intended interpretation is via unicellular maps:
\[
  \defp(\p)=2g(\p).
\]
The map model is the standard two-permutation encoding of a one-face map: the
long cycle records the face order and the permutation \(\p\) records the
vertex identifications.  Euler's formula then converts the cycle-count defect
into twice the genus.  The present article does not prove the full map-model
equivalence; it isolates it as part of the Chapuy input below.
The target recurrence is a Chapuy slicing estimate of the form
\[
  N_{m,g}\leq (2n)^3 N_{m,g-1},
  \qquad n=m+1,
\]
where
\[
  N_{m,g}=\#\{\p:\defp(\p)=2g\}.
\]
Together with the genus-zero base, this gives
\[
  N_{m,g}\leq 4^n(2n)^{3g}.
\]
This is the enumerative heart of the remaining U-AUB-02 proof.

\section{The Biane Contour Mechanism}

The genus-zero injection depends on recovering a left-tight permutation from
its cycle partition.  This recovery is obstructed by predecessor skips.  The
proof route uses Biane's noncrossing geometry: from a skip one extracts a long
interval whose contour has a repeated class; path uniqueness in the associated
incidence tree then forces a contradiction.

Let
\[
  S_\p=\text{cycle partition of }\g\p,\qquad
  T_\p=\text{cycle partition of }\p.
\]
The incidence graph has left vertices the \(S_\p\)-classes, right vertices the
\(T_\p\)-classes, and an edge for each point of \(\{0,\ldots,m\}\).  In the
left-tight noncrossing case this incidence graph is a tree.  Hence two simple
paths with the same endpoints must coincide.

\begin{definition}[Contour over an interval]
For an interval \([a,c]\subseteq\{0,\ldots,m\}\), the Biane contour is the
alternating path in the incidence graph obtained by following the cyclic
successor steps
\[
  a\to a+1\to \cdots\to c.
\]
At each step the path crosses from an \(S_\p\)-class to a \(T_\p\)-class and
back.  A repeated contour class means that two distinct points of the interval
give the same left vertex \(S_\p\), or two distinct interior points give the
same right vertex \(T_\p\).  In the incidence tree, such a repeat is the
combinatorial event used to manufacture a second simple path.
\end{definition}

\begin{definition}[Predecessor skip]
A non-fixed point \(x\) is a predecessor skip if the cyclic predecessor
relation expected from the ordered cycle partition fails at \(x\).  Concretely
there are two shapes.
\begin{enumerate}[label=(\roman*)]
\item Nonwrapping: \(\p x<x\), and the \(\p\)-cycle of \(x\) contains a point
strictly between \(\p x\) and \(x\).
\item Wrapping: \(\p x\not <x\), and the \(\p\)-cycle of \(x\) contains a point
strictly below \(x\) or strictly above \(\p x\).
\end{enumerate}
\end{definition}

The direct-hit case is already elementary: if \(x\) and \(\g(\p x)\) lie in
the same \(\p\)-cycle, then one immediately obtains two distinct points lying
simultaneously in one \(T_\p\)-class and one \(S_\p\)-class.  This contradicts
the tree transversality of a left-tight noncrossing permutation.

The remaining case is the genuinely nested case:
\[
  x\not\sim_\p \g(\p x).
\]
The following lemma records the normalized interval data now available in all
three skip shapes.

\begin{lemma}[Branch-normalized nested return interval]
Let \(\p\) be left-tight and let \(x\) be a non-fixed predecessor skip.  Assume
the direct-hit case is excluded:
\[
  x\not\sim_\p \g(\p x).
\]
Then there exist \(a<c\) such that
\[
  c-a>1,\qquad a\sim_\p c,
\]
and the Biane contour over \([a,c]\) has a repeated \(S_\p\)-class or a
repeated \(T_\p\)-class.  Moreover the interval is canonical in its branch:
\begin{enumerate}[label=(\roman*)]
\item in the nonwrapping branch, \(a=\p x\), and \(c\) is the first same-\(\p\)
return after \(\p x\);
\item in the wrap-below branch, \(c=\p x\), and \(a\) is the last same-\(\p\)
return below \(x\);
\item in the wrap-above branch, \(a=x\), and \(c\) is the first same-\(\p\)
return above \(\p x\).
\end{enumerate}
\end{lemma}

\begin{proof}
In the nonwrapping branch, take the least same-cycle point \(c\) satisfying
\[
  \p x<c<x.
\]
If \(c=\p x+1\), then \(c=\g(\p x)\), contradicting the excluded direct-hit
case.  Hence \(c-\p x>1\).  Left-tightness then applies to the interval
\([\p x,c]\) and gives the repeated-contour obstruction.

In the wrap-below branch, nonfixedness and \(\p x\not<x\) imply
\[
  x<\p x.
\]
Take the largest same-cycle point \(a<x\).  Then \(\p x-a>1\), and the
interval \([a,\p x]\) carries the same repeated-contour obstruction.

In the wrap-above branch, take the least same-cycle point \(c>\p x\).  As in
the nonwrapping case, the direct-hit exclusion rules out \(c=\p x+1\).  Since
\(x<\p x<c\), the interval \([x,c]\) has length \(>1\), and left-tightness
again supplies the repeated-contour obstruction.
\end{proof}

\section{Finite Gap-Three Biane Reductions}

The adjacent-link theorem is the local Biane statement used in the count.  In
one orientation, it has the following form.  Suppose
\[
  a<b<c<d,\qquad a\sim_\p c,\qquad b\sim_\p d,
  \qquad a\not\sim_\p b.
\]
The desired left adjacent link is
\[
  a\sim_{S_\p} b.
\]
The gap-two case is already settled.  The gap-three cases split into small
finite contour branches.  The current reduction discharges these finite
branches by incidence-tree uniqueness.

The branch accounting used in the reduction is as follows.  Here \(S(u,v)\)
means \(u\sim_{S_\p}v\), and \(T(u,v)\) means \(u\sim_\p v\).

\begin{center}
\begin{tabular}{L{0.28\linewidth}L{0.36\linewidth}L{0.25\linewidth}}
\hline
Slice & Branches exposed by the contour repeat & Status in the reduction\\
\hline
\(c-a=3,\ b=a+1\)
&
\(S(a,a+2)\) with \(\neg T(a,a+2)\);
\(S(b,c)\);
fixed \(a+2\);
\(T(b,a+2)\)
&
contradicted or routed by tree uniqueness and the right frontier\\
\(c-a=3,\ b=a+2\)
&
fixed \(a+1\);
\(S(a+1,c)\);
\(T(a+1,b)\);
adjacent \(S(a+1,b)\)
&
absorbed by path uniqueness and adjacent-repeat normalization\\
\(d-b=3,\ c=b+2\)
&
\(S(b,c)\), \(S(b+1,c)\), \(S(b+1,d)\),
\(T(b,b+1)\), \(T(b,c)\), \(T(b+1,c)\)
&
absorbed by the right special-slice reductions\\
\hline
\end{tabular}
\end{center}

For example, in the slice
\[
  c-a=3,\qquad b=a+1,
\]
the remaining branch
\[
  a\not\sim_\p a+2,\qquad a\sim_{S_\p} a+2
\]
is contradictory.  The repeat \(a\sim_{S_\p} a+2\) copies one contour step
into a path whose endpoints are already connected by the direct path coming
from \(a\sim_\p c\).  Tree uniqueness forces \(a\sim_\p a+2\), contradiction.

In the other left gap-three slice,
\[
  c-a=3,\qquad b=a+2,
\]
the fixed branch at \(a+1\) and the branch \(a+1\sim_{S_\p}c\) are similarly
absorbed by comparing two simple incidence paths.  The checked right special
slice
\[
  d-b=3,\qquad c=b+2
\]
is handled by the same principle, with left and right contour roles shifted.

\section{The Conditional Bidefect Theorem}

We now state the current reduction in journal-readable form.  The hypotheses
are exactly the remaining theorem-strength leaves.

\begin{hypothesis}[Large-gap left adjacent links]
For every left-tight \(\p\) and every crossing quadruple \(a<b<c<d\) with
\[
  a\sim_\p c,\qquad b\sim_\p d,\qquad a\not\sim_\p b,
\]
if
\[
  2<c-a,\qquad c-a\neq 3,
\]
then the left adjacent link holds:
\[
  a\sim_{S_\p} b.
\]
\end{hypothesis}

\begin{hypothesis}[Large-gap right adjacent links]
Under the same crossing assumptions, if
\[
  2<d-b
\]
and the special checked slice
\[
  d-b=3,\qquad c=b+2
\]
does not occur, then the right adjacent link holds:
\[
  c\sim_{S_\p} d.
\]
\end{hypothesis}

\begin{hypothesis}[Return-normalized nested skip collision]
Let \(\p\) be left-tight and assume its cycle partition is noncrossing.  If
\(x\) is a non-fixed predecessor skip and
\[
  x\not\sim_\p \g(\p x),
\]
then the branch-normalized return interval supplied by the preceding lemma
forces two distinct points \(y,z\) such that
\[
  y\sim_\p z
  \qquad\text{and}\qquad
  y\sim_{S_\p}z.
\]
\end{hypothesis}

\begin{hypothesis}[Chapuy slicing bound]
The plus-defect classes admit the unicellular-map genus interpretation
\(\defp=2g\), and the corresponding one-sided genus slices satisfy
\[
  N_{m,g}\leq (2(m+1))^3N_{m,g-1}
\]
for all \(g\geq 1\), with the genus-zero base bounded by \(4^{m+1}\).
\end{hypothesis}

\begin{theorem}[Conditional bidefect Aubrun bound]
Assume the four hypotheses above.  Then for every \(m\geq 0\) and every
\[
  0\leq a\leq m+3,\qquad 0\leq b\leq m+1,
\]
one has
\[
  N_m(a,b)\leq \Env_m(a,b).
\]
\end{theorem}

\begin{proof}
The rank budget and bidefect split reduce the analytic Wick sum to a
two-parameter count over \((a,b)\).  The plus-defect class controls every
bidefect class with the same \(a\), so it remains to prove the one-sided
plus-defect bound.

The genus-zero plus-defect slice is the left-tight class.  Noncrossing and
predecessor reconstruction inject it into noncrossing partitions, giving the
bound \(4^{m+1}\).  The predecessor reconstruction follows by contradiction:
a non-fixed predecessor skip first gives either the direct collision, already
closed, or the nested collision supplied by the third hypothesis.  In either
case two distinct points lie in both a \(\p\)-cycle and an \(S_\p\)-cycle,
contradicting left-tight transversality.

The finite gap-three Biane branches used in this reconstruction are accounted
for in the table of Section~5 and handled by incidence-tree path-uniqueness
reductions.  The remaining large-gap adjacent links are exactly the first two
hypotheses.

The positive plus-defect slices are then bounded by the Chapuy recurrence.
Iterating from the genus-zero base gives the required polynomial envelope in
the plus defect.  Finally, the bidefect scalar absorption attaches the two
binomial factors and pays the two defects by separate shifted bases.  This is
precisely the envelope \(\Env_m(a,b)\).
\end{proof}

\section{What Remains for an Unconditional Article}

The theorem is a clean reduction, not yet the final Aubrun count theorem.  To
remove the word ``conditional'', one must prove the four hypotheses in
Section~6.

The first two are large-gap Biane adjacent-link theorems.  The finite
gap-three slices are no longer the issue; the remaining cases require a
uniform argument for long contour gaps.  The third is the nested skip
collision theorem.  The branch-normalized return interval lemma reduces its
bookkeeping: every branch now provides a canonical long interval with a
repeated contour class, and the current formal reduction makes precisely this
return interval theorem-facing.  The fourth is the Chapuy recurrence together
with the exact conversion of plus-defect into unicellular genus.

\appendix
\section{Synchronization with the Formal Development}
\label{sec:sync}

This appendix is not part of the mathematical proof.  It records the current
source of the statements used above.

\begin{description}[font=\normalfont\bfseries]
\item[Bidefect split]
\path{wickPlusDefect_add_wickMinusDefect_eq_wickDefect}
\item[Bidefect envelope]
\path{aubrunBiDefectCountEnvelope} and
\path{AubrunBiDefectCountBound}
\item[Plus-defect to bidefect count]
\path{wickBiDefectCount_le_biDefectEnvelope_of_plusDefectCount}
\item[Genus-zero left-tight base]
\path{wickPlusGenusCount_zero_le_four_pow_of_no_nonfixed_cyclicPredecessorSkip}
\item[Direct predecessor-skip collision]
\path{permCyclicPredecessorSkip_direct_commonCycle}
\item[Branch-normalized nested interval]
\path{leftTight_nonfixed_nested_cyclicPredecessorSkip_exists_returnContourIntervalRepeat}
\item[Current conditional endpoint]
\path{AubrunBiDefectCountBound_of_bianeAdjacentLinkHalves_left_gap_three_start_impossible_frontiers_nestedReturnContourIntervalCommonCycle_and_chapuy}
\end{description}

The current audit of the return-normalized endpoint reports only the
standard foundational dependencies \(\texttt{propext}\),
\(\texttt{Classical.choice}\), and \(\texttt{Quot.sound}\).  This is not a
claim that U-AUB-02 is done; it is evidence that the nested collision leaf is
now exposed through the checked branch-normalized return interval.

\begin{thebibliography}{9}

\bibitem{Aubrun}
G. Aubrun,
\emph{Partial transposition of random states and non-centered semicircular
distributions},
Random Matrices: Theory and Applications \textbf{1} (2012), no. 2, 1250001.

\bibitem{Biane}
P. Biane,
\emph{Some properties of crossings and partitions},
Discrete Mathematics \textbf{175} (1997), no. 1--3, 41--53.

\bibitem{Chapuy}
G. Chapuy,
\emph{A new combinatorial identity for unicellular maps, via a direct
bijective approach},
Advances in Applied Mathematics \textbf{47} (2011), no. 4, 874--893.

\bibitem{HarerZagier}
J. Harer and D. Zagier,
\emph{The Euler characteristic of the moduli space of curves},
Inventiones Mathematicae \textbf{85} (1986), no. 3, 457--485.

\end{thebibliography}

\end{document}
